%! Author = sriadityavedantam
%! Date = 3/29/26

\section{Relations}

\begin{definition}[Relations and Cartesian Products]
    A \textbf{cartesian product} of sets $A$ and $B$ gives us
    \[
        A \times B \coloneqq \{(a,b) : a \in A,\ b \in B\}.
    \]
    Note that
    \[
        A \times \varnothing = \varnothing = \varnothing \times B.
    \]
    An equivalence relation, denoted by $\sim$, on a set $A$ is a relation satisfying the following properties:
    \begin{itemize}
        \item \textbf{reflexive:} $a \sim a$ for all $a \in A$;
        \item \textbf{symmetric:} if $a \sim b$, then $b \sim a$ for all $a,b \in A$;
        \item \textbf{transitive:} if $a \sim b$ and $b \sim c$, then $a \sim c$ for all $a,b,c \in A$.
    \end{itemize}
\end{definition}

\section{Well Ordering and Induction}

\begin{definition}[Well-Ordering]
    Every nonempty subset of $\mathbb{Z}^{\geq 0}$ contains a smallest element.
\end{definition}

This takes into account that there is an order relation $(<)$ on all integers of $\mathbb{Z}$.
The direct consequence of this definition is mathematical induction.
Mathematical induction is a proof technique that uses recursive techniques to prove that a statement is true for all elements past its base case.

\begin{theorem}{Assume that $n \in \mathbb{Z}^{\geq 0}$ and $P(n)$ is given.
    \begin{enumerate}
        \item $P(0)$ is a true statement.
        \item When $P(k)$ is true, then $P(k+1)$ is also true.
    \end{enumerate}
    Then $P(n)$ is true for all $n \in \mathbb{Z}^{\geq 0}$.
}
A remark on this theorem is that $P(k)$ does not have to be true, but we assume so.
This is called the induction hypothesis.

In proof writing, if we are given an ``If\ldots Then\ldots'' statement, we generally assume that the statement before the ``Then'' is true, and attempt to prove the rest.
This is the same thing we have proved through induction.
It can be seen as a result of continued direct proofs compiled together and generalized to become the induction we know today.

The following example is how we use induction in today's world, and it is important to note how we use it compared to how one may have done it for a proofs course.
In other words, this is a practical application of how a researcher would use induction.
\end{theorem}

\begin{example}
    A set of $n$ elements has $2^n$ subsets.

    $P(0)$:
    $2^0 = 1$ subset.

    $P(1)$:
    $2^1 = 2$ subsets.

    $P(3)$:
    $2^3 = 8$ subsets.

    Assume $P(k)$, namely that a set with $k$ elements has $2^k$ subsets.
    Now prove that $P(k+1)$ says a set with $k+1$ elements has $2^{k+1}$ subsets.
    In a more standardized proof writing, we can define a set
    \[
        S \coloneqq \{n \in \mathbb{Z}^{\geq 0} : P(n) \text{ is true}\},
    \]
    and show that $S = \mathbb{Z}^{\geq 0}$.
    Let our induction hypothesis be ``$P(k)$ is true.''
    Since we have shown that our base case $P(0)$ is true, we assume $P(k)$ is true and attempt to prove $P(k+1)$.
    Suppose a set has $k$ elements.
    If we add one new element, then every old subset has exactly two choices:
    either include the new element or do not include it.
    Therefore the number of subsets doubles.
    Hence the new set has
    \[
        2 \cdot 2^k = 2^{k+1}
    \]
    subsets.
    Thus $P(k+1)$ is true.
    Therefore, by induction, $P(n)$ is true for all $n \in \mathbb{Z}^{\geq 0}$.
\end{example}